% Part IV -- 3D-stacked cache and AMD 3D V-Cache (target: about 25 pages)
\clearpage
\part{3D-Stacked Cache and AMD 3D V-Cache}
\label{part:stacked-cache}
\partbudget{25}{argument for stacking, planar limits, side-by-side versus vertical comparison and hybrid-bonding facts drafted; thermal and workload chapters need independent sources}

\section{Why Cache Capacity Matters}
\label{sec:why-cache-capacity}

\Cref{sec:amat,sec:miss-classes} showed that a larger LLC lowers the miss rate by removing capacity misses, provided its hit time does not grow too much. Whether that matters depends on the workload: a program whose working set already fits in the cache gains little, while one whose working set lies just above the cache size can gain a great deal.

AMD's ISSCC~2022 presentation of 3D V-Cache opens with two charts that frame this part \cite[slides~5--6]{wuu2022vcache}. The first, titled ``Large L3 Caches Provide IPC Uplift'', plots IPC uplift on a linear scale rising steadily as L3 size grows from 1$\times$ to 32$\times$; the slide gives no values on the uplift axis and does not name the workloads. The second, ``Silicon Area Scaling by Function'', shows logic area continuing to shrink from 28\,nm to 5\,nm while SRAM scaling falls behind from about 7\,nm and analog circuits shrink least; a companion chart on the same slide shows normalised cost per yielded mm$^2$ (for a 250\,mm$^2$ die) rising from 45\,nm to 5\,nm. The slide summarises these as ``SRAM scaling \& cost = barriers to large on-die caches''. This extends Naffziger \emph{et al.}'s 2020 observation that analog circuits scale worst \cite{naffziger2020chiplet}: by the 5\,nm generation SRAM, too, has become comparatively expensive.

Both charts are vendor material, and the first has no stated workloads or values; neither is evidence that every program benefits. Part~\ref{part:case-study} tests the claim against published measurements.

\section{Why Expanding Cache in Two Dimensions Became Difficult}
\label{sec:planar-limits}

Growing the L3 in the plane of the compute die meets four related obstacles.

\subsection{Area and yield}
SRAM cells no longer shrink at the rate of logic (\cref{sec:why-cache-capacity}), so a cache of a given capacity occupies a growing share of the die at each new node. AMD lists cost, product flexibility and latency as the resulting barriers to large on-die caches \cite[slide~6]{wuu2022vcache}. Mittal and Vetter add that, if the large caches needed by future workloads were built from SRAM, they could occupy up to 90\% of the chip area \cite[sec.~1]{mittalSurveyTechniquesArchitecting2016}. A larger die yields fewer good dies per wafer (\cref{sec:yield-cost}), so doubling the L3 on a planar CCD would raise the cost of every CCD, including those sold for workloads that do not need the extra cache. It would also make the CCD less flexible as a building block reused across products.

\subsection{Wire length and access latency}
A larger planar cache spreads over a larger area, so the average distance between a core and the cache bank holding its data increases. Longer wires have more resistance and capacitance; their RC delay increases access latency, and the energy needed to charge their capacitance increases dynamic power (\cref{eq:pdyn}). In terms of \cref{eq:amat}, part of the miss-rate gain from a bigger cache is paid back as a longer hit time $t_{\mathrm{LLC}}$. AMD makes the same argument for stacking: compared with spreading a large cache across the plane, 3D integration ``improves effective memory latency'' and ``reduces long datapath and I/O's dynamic powers'' \cite[slide~7]{wuu2022vcache}.

\subsection{Power}
Large SRAM caches consume dynamic power on every read and write and leakage power even when idle (\cref{sec:power}). Mittal and Vetter's survey of DRAM-cache architectures argues that the low density and high leakage power of SRAM make very large SRAM caches infeasible, and that such caches may need expensive cooling, such as liquid cooling, which adds substantial cost and complexity \cite[sec.~1]{mittalSurveyTechniquesArchitecting2016}. This is one reason researchers have explored denser technologies, such as stacked DRAM, for last-level caches.

\subsection{Summary of the planar limits}
Planar expansion therefore trades miss rate against die cost, hit latency and power all at once. The question for the next section is whether a separate cache die can break that coupling.

\section{Side-by-Side Versus Vertically Stacked Cache Dies}
\label{sec:horizontal-vs-vertical}

If extra cache is to be made on its own die, it can be placed either \emph{beside} the compute die in the package or \emph{on top of} it (\cref{fig:horizontal-vs-vertical}). Both are valid chiplet designs; the choice is a trade-off rather than a matter of one being impossible.

\input{"../draft-initial/figures/horizontal-vs-vertical.tex"}

\subsection{Placing the cache die beside the compute die}
A side-by-side cache die is a perfectly workable chiplet design, but it has three costs:
\begin{itemize}
  \item \textbf{Package footprint.} The extra die occupies additional area next to the CCD. This enlarges the \emph{package}, even though it does not enlarge the compute die itself. The distinction between \emph{die area} (one piece of silicon) and \emph{package footprint} (the area of the assembled processor) is important throughout this part.
  \item \textbf{Interconnect.} Every access to the extra cache must cross a die-to-die interface through the package or an interposer. Package-level connections are far coarser than on-die wiring, which limits how many parallel connections---and hence how much bandwidth---the link can offer.
  \item \textbf{Latency and power.} Data must travel between the dies, adding interface circuitry and wire delay to every access to the extra cache, and energy to every bit moved.
\end{itemize}

A side-by-side die does not automatically push the CCD farther from the IOD: the package designer chooses where each die goes, and the cache die could be placed on the far side of the CCD. However, package area is limited, the CCD--IOD links are already routed with care to minimise memory latency \cite{naffziger2020chiplet}, and every additional die competes for the same routing resources. At high signalling rates, longer and more numerous package routes bring RC delay, power, signal-integrity and routing-congestion problems. These can be mitigated with better interconnect design, buffering and packaging technology, so they are engineering costs rather than impossibilities.

The more important comparison is therefore not the CCD-to-IOD distance but the latency and bandwidth of the path \emph{between the cores and the additional cache}.

AMD's own framing of the choice matches this analysis. Its ISSCC~2022 presentation illustrates a hypothetical processor whose large cache is split off onto a side-by-side chiplet, and states that ``2.5D chiplets can provide product flexibility and reduce cost'' but that ``3D can be even better'', because it improves effective memory latency, reduces long-datapath and I/O dynamic power, and ``fits more transistors within a given package cavity size'' \cite[slide~7]{wuu2022vcache}.

\subsection{Stacking the cache die on the compute die}
Vertical stacking changes the trade-off:
\begin{itemize}
  \item \textbf{Footprint.} The cache die overlaps the CCD's existing footprint, so capacity grows with little additional package area.
  \item \textbf{Interconnect density.} The two dies can be joined by a very large number of short vertical connections, which supports the bandwidth a large cache needs and keeps the added distance small.
  \item \textbf{Thermal design.} The stack can obstruct heat flow from the compute die, depending on which die is on top and what lies above the cores. This is the central new engineering problem (\cref{sec:thermal}).
\end{itemize}

\begin{table}[htbp]
  \centering
  \caption{Qualitative comparison of side-by-side and vertically stacked cache dies}
  \label{tab:horizontal-vs-vertical}
  \small
  \begin{tabularx}{\textwidth}{@{}>{\raggedright\arraybackslash}p{3.1cm}>{\raggedright\arraybackslash}X>{\raggedright\arraybackslash}X@{}}
    \toprule
    \textbf{Factor} & \textbf{Cache die beside the CCD} & \textbf{Cache die stacked on the CCD} \\
    \midrule
    Package footprint & Larger & More compact \\
    Connections & Side-by-side package or interposer link & Dense vertical bonds and through-silicon vias \\
    Cache access latency & Depends on the interface design & Can benefit from short, dense connections \\
    Cooling & Easier to arrange & More difficult; depends on stack order \\
    Manufacturing & Extra die plus package assembly & Extra die plus bonding, thinning and stacking \\
    \bottomrule
  \end{tabularx}
\end{table}

\Cref{tab:horizontal-vs-vertical} summarises the comparison. The case for vertical stacking is not that side-by-side chiplets cannot work. It is that stacking can add substantial cache capacity without enlarging the package by the same amount, with potentially better interconnect characteristics. The question worth investigating is \emph{when the extra complexity of 3D integration is worth the smaller footprint and denser interconnect}.

\section{Stacked-Die Integration and Hybrid Bonding}
\label{sec:hybrid-bonding}

\subsection{The first-generation 3D V-Cache stack}
AMD's Hot Chips~33 presentation describes the first-generation stack as a 64\,MB L3 cache die attached to an up-to-eight-core Zen~3 CCD by a direct copper-to-copper bond, with \glspl{tsv} carrying signals and power through the CCD \cite{swaminathanAdvancedPackagingEnabling2021}. AMD's EPYC~7003 documentation states that the L3 cache tiles are stacked on the central L3 portion of the base CCX \cite{amd2022epyc7003}. The ISSCC~2022 presentation describes the three components of the stack and how they are assembled \cite[slides~10--12, 14, 16--17]{wuu2022vcache}:
\begin{itemize}
  \item the Zen~3 CCD, built on TSMC's 7\,nm process with eight cores sharing 32\,MB of L3 in an area of 81\,mm$^2$, with V-Cache support designed in ``from Day~1'';
  \item the extended L3 die (L3D), also 7\,nm, carrying 64\,MB of L3 in 41\,mm$^2$, roughly half the CCD area;
  \item structural silicon dies, which provide mechanical support for the thinned CCD and a path for heat from the CPU cores.
\end{itemize}
The CCD is mounted face-down on the package substrate through conventional C4 bumps and connects to the L3D through TSVs. The L3D is mounted face-down on the CCD by hybrid bonding, using TSMC's SoIC process with a minimum TSV pitch of 9\,\textmu m, while the structural dies are oxide-bonded beside it. Support silicon brings the stack to the same height as an unstacked CCD, so server parts with and without V-Cache fit the same package. The cache die sits over the CCD's central L3 region \cite{amd2022epyc7003} and the structural silicon over the cores, so the hottest circuits are covered only by passive silicon.

\subsection{Hybrid bonding versus micro-bumps}
Earlier 3D stacking typically joined dies with solder \glspl{micro-bump}. \Gls{hybrid-bonding} instead bonds aligned copper pads and surrounding dielectric directly. AMD reports, from its internal engineering analysis \cite{swaminathanAdvancedPackagingEnabling2021}:\index{hybrid bonding}
\begin{itemize}
  \item a hybrid-bond pitch of 9\,\textmu m, compared with 50 to 36\,\textmu m for micro-bump 3D;
  \item more than 15$\times$ the interconnect density of micro-bump 3D;
  \item more than 3$\times$ better interconnect power efficiency than micro-bump 3D;
  \item lower TSV capacitance and inductance, which benefits signal and power delivery.
\end{itemize}
AMD's EPYC~7003 documentation similarly claims over 200$\times$ the interconnect density of 2D packaging and over 15$\times$ that of solder-bump 3D \cite{amd2022epyc7003}. These are AMD's own comparisons and should be read as engineering claims.

\plannedcontent{Explain why interconnect density matters architecturally (parallel connections $\rightarrow$ bandwidth; short connections $\rightarrow$ low latency and energy), and how this addresses the side-by-side costs in \cref{sec:horizontal-vs-vertical}. Keep fabrication detail at the architectural level.}

\section{Cache Access, Latency, Bandwidth and Hierarchy}
\label{sec:vcache-access}

With the stacked die, the L3 of a Zen~3 CCD grows from 32\,MB to 96\,MB, shared by up to eight cores \cite{wuu2022vcache,amd2022epyc7003}; an EPYC~7003 processor with eight such CCDs carries up to 768\,MB of L3 \cite{amd2022epyc7003}. AMD reports up to 2\,TB/s of bandwidth for the V-Cache \cite{wuu2022vcache}. The extra capacity remains part of the same L3 level, so the cache hierarchy above it (private L1 and L2 per core) is unchanged \cite{amd2022epyc7003}.

AMD's presentation gives the cost of the extra capacity in latency: the 96\,MB L3 remains 16-way set associative with a 32\,B/cycle interface to each core, delivers more than 2\,TB/s of L3 bandwidth, and adds \emph{four cycles} of latency \cite[slide~20]{wuu2022vcache}. Each die's portion of the L3 has its own data, tag and LRU arrays, so a tag lookup does not have to cross between dies before the data array is accessed \cite[slide~20]{wuu2022vcache}.

\Cref{eq:amat} shows how to judge this trade-off. For an access that has already missed in L2, the extra capacity changes the average cost by approximately
\begin{equation}
  \Delta T \;\approx\; \Delta t_{\mathrm{L3}} \;-\; \Delta m_{\mathrm{L3}} \cdot t_{\mathrm{mem}},
  \label{eq:vcache-tradeoff}
\end{equation}
where $\Delta t_{\mathrm{L3}} = 4$ cycles is the added hit latency, $\Delta m_{\mathrm{L3}}$ the reduction in the L3 miss rate and $t_{\mathrm{mem}}$ the main-memory latency in cycles. As a purely illustrative example, if main memory took 300 cycles, the larger L3 would pay off once it removed more than about $4/300 \approx 1.3$ percentage points of L3 misses. A workload whose working set already fits in 32\,MB gains nothing and pays the four cycles on every L3 hit; a workload whose working set lies between 32\,MB and 96\,MB can gain substantially. \todo{Replace the illustrative 300-cycle memory latency with a cited value, or keep it explicitly labelled as hypothetical.}

The L3 remains private to each CCD, so a thread running on one CCD does not benefit from another CCD's stacked cache.

\section{Thermal Management, Power Delivery and Manufacturing Complexity}
\label{sec:thermal}

Stacking changes the heat path described in \cref{sec:power}: heat from the cores must now pass through additional silicon and bonding layers before it reaches the heat spreader. The first-generation design limits the penalty by placing the cache die over the CCD's central L3 region \cite{amd2022epyc7003} and covering the cores only with structural silicon, whose stated purposes include ``thermal dissipation for CPU cores'' \cite[slide~12]{wuu2022vcache}. AMD also credits hybrid bonding with ``superior thermal conductance'' compared with micro-bump stacking \cite[slide~15]{wuu2022vcache}, although this is an AMD claim based on internal simulation. Earlier academic work points the same way: in a study reviewed by Mittal and Vetter, Black \emph{et al.} stacked an 8\,MB SRAM cache directly on a dual-core processor and found that none of their stacked-cache designs changed the thermal profile significantly, although stacked SRAM raised temperature more than stacked DRAM because of its higher power density \cite[p.~1853]{mittalSurveyTechniquesArchitecting2016}.

Power must also be delivered upwards. The L3D receives its supplies through power TSVs placed in channels between the CCD's array macros, and its arrays use separate rails for the bit cells and the peripheral circuits \cite[slides~24--25]{wuu2022vcache}.

\plannedcontent{Die thinning, bonding and test complexity (design-for-test is listed among the topics of \cite{wuu2022vcache}); frequency and voltage limits of stacked-cache parts. \todo{Identify independent sources on thermal behaviour and on frequency or voltage limits of stacked-cache parts.} \verify{Later generations are reported to place the cache die beneath the CCD; confirm from a primary source before mentioning.}}

\section{Workloads That Benefit From a Larger LLC}
\label{sec:vcache-workloads}

AMD reports a 15\% average gaming performance uplift from V-Cache \cite{wuu2022vcache}. Its Hot Chips~33 comparison of a 12-core 3D-chiplet prototype with a Ryzen~9 5900X at the same fixed frequency reports an average gain of 15\% across a small set of games \cite{swaminathanAdvancedPackagingEnabling2021}. Holding frequency fixed is valuable because it isolates the cache effect, but the result covers only the disclosed games and configuration and should not be generalised.

The ISSCC~2022 presentation adds product results, again from AMD's own testing \cite[slides~29, 31, 35]{wuu2022vcache}. The eight-core Ryzen~7 5800X3D is reported as about 15\% faster in games than the twelve-core Ryzen~9 5900X at 1080p with high settings, ranging from up to 1.36$\times$ in one title to a tie in another. On servers, a 16-core EPYC with V-Cache completed about 66\% more electronic-design (RTL simulation) jobs per hour than a 16-core EPYC~73F3 without it. The spread from a tie to 1.36$\times$ already illustrates the workload dependence predicted in \cref{sec:miss-classes}; \cref{sec:confounds} discusses what else differs between the compared parts.

\plannedcontent{Characterise cache-sensitive workloads in terms of working-set size and reuse (\cref{sec:miss-classes}): games, some simulation and electronic-design workloads, databases. Characterise insensitive workloads (small working sets; streaming data with little reuse). All quantitative claims to come from Part~\ref{part:case-study}.}

\section{Summary of Part IV}
\plannedcontent{Three-part test: \emph{bottleneck} (planar cache cost, latency and power); \emph{benefit} (tripled L3 per CCD with dense, short connections); \emph{new trade-offs} (heat removal, bonding and test complexity, product cost, workload dependence).}
